% =============================================================================
% style/base_subsection.tex
% =============================================================================


% =============================================================================
% Цвета заголовков подразделов
% =============================================================================

\colorlet{subsectionAccent}{xred!72!black}
\colorlet{subsectionTitleColor}{subsectionAccent}

% Три тонкие вертикальные линии (subsection всегда без номера).
\colorlet{subsectionMarkA}{xred!60!white}
\colorlet{subsectionMarkB}{xred!75!white}
\colorlet{subsectionMarkC}{xred!90!black}


% =============================================================================
% Геометрия
% =============================================================================

\newlength{\subsectionBaselineShift}
\setlength{\subsectionBaselineShift}{4.6pt}


% =============================================================================
% Метка subsection: три тонкие вертикальные линии
% =============================================================================

\providecommand{\subsectionMark}{}
\renewcommand{\subsectionMark}{%
  \tikz[baseline={([yshift=-\subsectionBaselineShift]S.center)}]{%
    \node[
      minimum width=5.2mm,
      minimum height=5.4mm,
      inner sep=0pt,
      outer sep=0pt
    ] (S) {};

    \draw[
      subsectionMarkA,
      line width=0.95pt,
      line cap=round,
    ]
      ([xshift=0.8mm, yshift=-0.15mm]S.north west) --
      ([xshift=0.8mm, yshift=0.15mm]S.south west);

    \draw[
      subsectionMarkB,
      line width=1.10pt,
      line cap=round,
    ]
      ([xshift=2.2mm, yshift=-0.15mm]S.north west) --
      ([xshift=2.2mm, yshift=0.15mm]S.south west);

    \draw[
      subsectionMarkC,
      line width=1.25pt,
      line cap=round,
    ]
      ([xshift=3.6mm, yshift=-0.15mm]S.north west) --
      ([xshift=3.6mm, yshift=0.15mm]S.south west);
  }%
}


% =============================================================================
% Subsection: метка слева, без горизонтальных линий
% =============================================================================

\titleformat{\subsection}[hang]
  {%
    \normalfont\sffamily\bfseries\fontsize{15.5}{19}\selectfont%
    \raggedright\color{subsectionTitleColor}%
    \hyphenpenalty=10000\exhyphenpenalty=10000%
  }
  {\subsectionMark}
  {0.55em}
  {}


% =============================================================================
% Звёздная форма subsection
% =============================================================================

\titleformat{name=\subsection,numberless}[hang]
  {%
    \normalfont\sffamily\bfseries\fontsize{15.5}{19}\selectfont%
    \raggedright\color{subsectionTitleColor}%
    \hyphenpenalty=10000\exhyphenpenalty=10000%
  }
  {\subsectionMark}
  {0.55em}
  {}


% =============================================================================
% Отступы subsection
% =============================================================================

\titlespacing*{\subsection}
  {0pt}
  {4.5ex plus 1.5ex minus 1.3ex}
  {2.0ex plus 0.5ex minus 0.4ex}